\begin{description}
\item[(a)] $L\propto P^{2.5}$ implies $L=KP^{2.5}$, so
\[\log L=\log K+2.5\log P.\]
\hfill(2 points)
Also $\log L=\log F+\log4+\log\pi+2\log r$.
The magnitude relation is $m_1-m_2=-2.5\log(F_1/F_2)$. \hfill(1 point)
For ease of calculation use Vega as reference star and let $m$ be the apparent magnitude of the Cepheid variable. Then $-m/2.5=\log F_1-\log F_0$. \hfill(1 point)
Combining this with $m=M+5\log(r/(10\,\mathrm{pc}))$ gives
\[M=-6.25\log P+C,\]
where $C=-2.5(\log K-\log4-\log F_0-\log\pi)+5\log10$. \hfill(3 points)
\item[(b)] The magnitude--luminosity relation has $C=-1.5$. \hfill(1 point)
Suppose $m$ is the average magnitude corresponding to the luminosity $\tfrac12(L_{\max}+L_{\min})$. \hfill(1 point)
Then
\[m-m_{\mathrm{ref}}=-2.5\log\left(\frac{(L_1+L_2)/2}{L_{\mathrm{ref}}}\right).\]
\hfill(1 point)
Equivalently,
\[m=m_{\mathrm{ref}}+2.5\log2-2.5\log\left(\frac{L_1+L_2}{L_{\mathrm{ref}}}\right).\]
Since
\[\frac{L_{\mathrm{tot}}}{L_{\mathrm{ref}}}=\frac{L_1}{L_{\mathrm{ref}}}+\frac{L_2}{L_{\mathrm{ref}}}=(10^{-0.4m_1}+10^{-0.4m_2})10^{0.4m_{\mathrm{ref}}},\]
\hfill(1 point)
substitution gives
\[m=2.5\log2-2.5\log(10^{-2m_1/5}+10^{-2m_2/5}).\]
\hfill(1 point)
From the graph $m_{\max}=28$, $m_{\min}=22$, so $m=22.75$, and the period is 20 days. \hfill(1 point)
The absolute magnitude relation gives $M=-9.63$. \hfill(1 point)
Substitute into $m=M+5\log(r/(10\,\mathrm{pc}))$ to obtain $r=29.92\,\mathrm{Mpc}$. \hfill(1 point)
\end{description}
